
\label{chapter: AppendixB}



The reason why the vectors $\boldsymbol{\tau}_\alpha$ have to be used for the Fourier transformation in Eq.~\eqref{eq: fourier transf} is that the Berry connection, and therefore the Berry curvature and Berry phase, require cell-periodic states $\ket{u_{n\mathbf{k}}}$ because the overlap of two Bloch functions
\begin{align}
   \langle{{\Psi_{n\mathbf{k}}}}|{{\Psi_{n\mathbf{k+b}}}}\rangle &= \int_{-\infty}^{\infty} \Psi_{n\mathbf{k}}^\dagger(x)\cdot \Psi_{n\mathbf{k+b}}(x) \,\text{d}x \\
   &=  \bra{{u_{n\mathbf{k}}}}e^{-i\mathbf{b}\cdot \mathbf{x}}\ket{{u_{n\mathbf{k+b}}}} = \int_{-\infty}^{\infty} e^{-i\mathbf{b}\cdot \mathbf{x}}u_{n\mathbf{k}}^\dagger(x)\cdot u_{n\mathbf{k+b}}(x) \,\text{d}\mathbf{x} = 0
\end{align}
vanishes, as discussed in detail in \cite{Vanderbilt_2018}. Hence, the resulting basis coefiicients, and hence the states have to be transformed, given by
\begin{align}
\ket{\Psi_{n\mathbf{k}}} &= \sum_\alpha v_{\alpha n}(\mathbf{k}) c_{\mathbf{k}\alpha}^\dagger \ket{0} \nonumber\\
&= \frac{1}{\sqrt{N}} \sum_{\mathbf{R},\alpha} 
v_{\alpha n}e^{i\mathbf{k}\cdot\mathbf{R}} \ket{\mathbf{R},\alpha}.
\end{align}
Here, $v_{\alpha n}$ describes the coefficients and with that the eigenstates of the Hamiltonian $\mathcal{H}(\mathbf{k})$. By again using Bloch's theorem, the cell-periodic states are determined to
\begin{align}
u_{n\mathbf{k}}(r) = e^{-i\mathbf{k}\mathbf{r}}\Psi_{n\mathbf{k}}(\mathbf{r}).
\end{align}
By adding the geometric position of the sublattices and selecting the sublattice $A$ as the origin, we obtain
\begin{align}
u_{n\mathbf{k}}(\mathbf{R}+\boldsymbol{\tau}_\alpha) = e^{-i \mathbf{k} \mathbf{R}}e^{-i\mathbf{k}\boldsymbol{\tau}_\alpha}\Psi_{n\mathbf{k}}(\mathbf{R}+\boldsymbol{\tau}_\alpha) = \frac{1}{\sqrt{N}}e^{-i\mathbf{k}\boldsymbol{\tau}_\alpha}v_{\alpha n}(\mathbf{k}).
\end{align}
Hence, the cell-periodic states are determined by
\begin{align}
 \ket{u_{n\mathbf{k}}} = \mathbf{D}(\mathbf{k})\cdot \ket{v_{n\mathbf{k}}}.
\end{align}
If, for example, the transformation is not used, all creation operators of one unit cell are mapped onto one vector, changing the tight-binding embedding used for geometric quantities, and therefore resulting in a different Berry curvature.
Lastly, the periodicity in $\mathbf{k}$ is shown for the Berry curvature
\begin{align}
    \Omega_n(\mathbf{k}+\mathbf{G}) &\approx \frac{1}{\text{d}A}\Im\Big\{\ln\big[
    \bra{v_{n\mathbf{k}+\mathbf{G}}} e^{-i\mathbf{G}\cdot\mathbf{r}} D( \Delta \mathbf{k}_x + \mathbf{G}) e^{i\mathbf{G}\cdot\mathbf{r}}\ket{v_{n\mathbf{k}+\mathbf{\Delta \mathbf{k}_x} + \mathbf{G}}}  \\
    &\quad+
    \bra{v_{n\mathbf{k}+\Delta \mathbf{k}_x + \mathbf{G}}} e^{-i\mathbf{G}\cdot\mathbf{r}} D( \Delta \mathbf{k}_y)e^{i\mathbf{G}\cdot\mathbf{r}}\ket{v_{n\mathbf{k}+\mathbf{\Delta \mathbf{k}_x} + \mathbf{k}_y + \mathbf{G}}} \\
    &\quad+
    \bra{v_{n\mathbf{k}+\Delta \mathbf{k}_x + \Delta \mathbf{k}_y + \mathbf{G}}} e^{-i\mathbf{G}\cdot\mathbf{r}}D(- \Delta \mathbf{k}_x)e^{i\mathbf{G}\cdot\mathbf{r}}\ket{v_{n\mathbf{k}+\Delta \mathbf{k}_y + \mathbf{G}}} \\
    &\quad+ 
    \bra{v_{n\mathbf{k}+\Delta \mathbf{k}_y + \mathbf{G}}} e^{-i\mathbf{G}\cdot\mathbf{r}} D( -\Delta \mathbf{k}_y) e^{i\mathbf{G}\cdot\mathbf{r}}\ket{v_{n\mathbf{k} + \mathbf{G}}}\big] \Big\}\\
    &= \frac{1}{\text{d}A}\Im\Big\{\ln\big[
    \braket{u_{n\mathbf{k}}}{u_{n\mathbf{k}+\mathbf{\Delta \mathbf{k}_x}}} +
    \braket{u_{n\mathbf{k}+\Delta \mathbf{k}_x}}{u_{n\mathbf{k}+\mathbf{\Delta \mathbf{k}_x}}} \\
    &\quad+
    \braket{u_{n\mathbf{k}+\Delta \mathbf{k}_x + \Delta \mathbf{k}_y}}{u_{n\mathbf{k}+\Delta \mathbf{\mathbf{k}_y}}} +
    \braket{u_{n\mathbf{k}+\Delta \mathbf{k}_y}}{u_{n\mathbf{k}}}\big] \Big\} \\
    &\approx \Omega_n(\mathbf{k}),
\end{align}
 where $u_{n\mathbf{k}}(r) = e^{i\mathbf{k}\mathbf{r}}v_{n,\mathbf{k}}(\mathbf{r})$ and $e^{-i\mathbf{G}\cdot\mathbf{r}}e^{i\mathbf{G}\cdot\mathbf{r}} = 1$ is used for the Fukui-hatsugai-Suzuki plaquette method.